\scanpage{47}
% Pencil archival foliation: 46.
d'après le n° précédent que \deleted{la composante \unsure{bipure} de}
\inserted{l'image dans}
\[
\varinjlim_{\substack{U\ \mathrm{ouvert}\\ \deleted{\mathrm{de}}\\ \mathrm{dense}}}
H^j(U,\unsure{\mathbb{Q}(0)|\mathbb{Z}(0)})
\]
de $H^j(X,\mathbb{Z}(0))$.

En fait, on constate (via dualité)
que si $X$ est lisse (sans doute
inutile) de dim $\leq n$, et si $i\leq n$
alors pour $\forall$ ouvert dense $U$ de $X$,
\[
H^i(X)\unsure{-|\longrightarrow}H^i(U)
\]
est \uline{\deleted{bijectif}\inserted{injectif} mod}
$D_{i-2}(\mathcal{M}^{+}(X))$ et
\uline{bijectif mod $D_{i-1}(\mathcal{M}^{+}(X))$}; ce qui
donne un invariant birationnel dans
\[
D_i(\mathcal{M}^{+}(X))/\unsure{D_{i-1}}(\mathcal{M}^{+}(X)),
\]
lequel
est d'ailleurs dans l'image de
$\mathcal{M}^{+}_{(i)}(X)$, donc définit un objet de
\[
\mathcal{M}^{+}_{(i)}(X)/
[\mathcal{M}^{+}_{(i-1)}(X)\unsure{+}\xi\mathcal{M}^{+}_{(i-2)}(X)],
\]
donc un
élément de $\mathbb{Z}^{\sigma_i(X)}$. C'est
l'invariant birationnel $\gamma_i(X)$.
